\section{{\LaTeX} examples}
\label{s:ex}

\paragraph{Commands.}
% use \sys
\sys helps students to use {\LaTeX} for writing a paper.
To comment on the paper, use
\TODO{to do.}
\XXX{fix this.}
\TK{it's mine.}

\paragraph{Pointing.}
% how to refer
Use \autoref{s:ex} to refer a section, 
use \autoref{t:simple} for a table,
and use \autoref{f:turtle} for a figure. 
To cite, use~\cite{kim:userfs,kim:poirot}. 
To refer a website, use \url{http://taesoo.org}. 

\paragraph{Forms.}

% compact list, [leftmargin=.2in]
\begin{itemize}[noitemsep,nolistsep]
 \setlength{\itemsep}{-0pt}
 \item foo
 \item bar
 \item baz
\end{itemize}

\paragraph{Generated figures.}

To render a svg, put the file to \cc{fig}, and include 
it to the latex like below. 

\begin{figure}[h]
\centering
\footnotesize
\includegraphics[width=0.5\columnwidth]{fig/ex-turtle}
\caption{A simple figure, generated from svg.}
\label{f:turtle}
\end{figure}

Similarly, to render a code, put a code snippet to \cc{code}, 
and include it to the latex like below.

\begin{figure}[h]
\centering
\footnotesize
\begin{Verbatim}[commandchars=\\\{\},codes={\catcode`\$=3\catcode`\^=7\catcode`\_=8}]
\PY{p}{[}\PY{c}{\PYZsh{}[\PYZsh{}[\PYZsh{}[\PYZsh{}[\PYZsh{}[\PYZsh{}[\PYZsh{}[\PYZsh{}[\PYZsh{} By TaroOgawa \PYZsh{}]\PYZsh{}]\PYZsh{}]\PYZsh{}]\PYZsh{}]\PYZsh{}]\PYZsh{}]\PYZsh{}]\PYZsh{}]}

                     \PY{n+nb}{globals}\PY{p}{(}\PY{p}{)}
               \PY{o}{.}\PY{n}{update}\PY{p}{(}\PY{p}{\PYZob{}}   \PY{l+s}{\PYZdq{}}\PY{l+s}{\PYZus{}\PYZus{}\PYZus{}\PYZus{}\PYZus{}\PYZus{}}\PY{l+s}{\PYZdq{}}\PY{p}{:}
            \PY{k}{lambda} \PY{n}{x}\PY{p}{:}         \PY{n+nb}{globals}\PY{p}{(}\PY{p}{)}
            \PY{o}{.}\PY{n}{update}\PY{p}{(}\PY{p}{(}         \PY{n+nb}{dict}\PY{p}{(}\PY{p}{[}\PY{p}{[}\PY{n}{x}\PY{p}{]}
               \PY{o}{*}\PY{l+m+mi}{2}\PY{p}{]}\PY{p}{)}\PY{p}{)}\PY{p}{)}\PY{p}{\PYZcb{}}\PY{p}{)}\PY{p}{,}   \PY{n}{\PYZus{}\PYZus{}\PYZus{}\PYZus{}\PYZus{}\PYZus{}}\PY{p}{(}\PY{p}{(}\PY{p}{(}
                     \PY{l+s}{\PYZdq{}}\PY{l+s}{Just}\PY{l+s}{\PYZdq{}}\PY{p}{)}\PY{p}{)}\PY{p}{)}
               \PY{p}{,}\PY{n}{\PYZus{}\PYZus{}\PYZus{}\PYZus{}\PYZus{}\PYZus{}}\PY{p}{(}\PY{p}{(}   \PY{l+s}{\PYZdq{}}\PY{l+s}{another}\PY{l+s}{\PYZdq{}}
            \PY{p}{)}\PY{p}{)}\PY{p}{,}\PY{n}{\PYZus{}\PYZus{}\PYZus{}\PYZus{}\PYZus{}\PYZus{}}         \PY{p}{(}\PY{l+s}{\PYZdq{}}\PY{l+s}{Python}\PY{l+s}{\PYZdq{}}
         \PY{p}{)}\PY{p}{,}\PY{n}{\PYZus{}\PYZus{}\PYZus{}\PYZus{}\PYZus{}\PYZus{}}\PY{p}{(}               \PY{l+s}{\PYZdq{}}\PY{l+s}{Hacker}\PY{l+s}{\PYZdq{}}\PY{p}{)}
         \PY{p}{]}\PY{p}{;}\PY{k}{print} \PY{p}{(}               \PY{l+s}{\PYZdq{}}\PY{l+s}{ }\PY{l+s}{\PYZdq{}}\PY{o}{.}\PY{n}{join}\PY{p}{(}
            \PY{p}{[}\PY{p}{(}\PY{n}{Just}\PY{p}{)}\PY{p}{,}\PY{p}{(}         \PY{p}{(}\PY{n}{another}\PY{p}{)}
               \PY{p}{)}\PY{p}{,}\PY{p}{(}\PY{n}{Python}   \PY{p}{)}\PY{p}{,}\PY{n}{Hacker}\PY{p}{]}
                     \PY{p}{)}\PY{p}{)}\PY{p}{;}\PY{n}{\PYZus{}\PYZus{}\PYZus{}\PYZus{}\PYZus{}\PYZus{}}
\end{Verbatim}

\caption{A highlighted code snippet, generated from 
 \cc{data/ex-code.py}
}
\label{f:code}
\end{figure}

Similarly, to render a \cc{gnuplot}, put a \cc{gnuplot}
script to \cc{data} with your actual data, 
and include it to the latex like below.

\begin{figure}[h]
\centering
\footnotesize
\includeplot{data/ex-gaussian}
\caption{A simple plot, generated directly from \cc{gnuplot}.}
\label{f:gaussian}
\end{figure}

\begin{figure}[h]
\centering
\footnotesize
\includeplot{data/ex-plot}
\caption{A simple plot, generated from dat.}
\label{f:plot}
\end{figure}

\paragraph{Table.}

% param: (cap-> force)
%  {h/H}ere
%  {t/T}op
%  {b/B}ottom
\begin{table}[h]
\centering
\footnotesize
\begin{tabular}{lcc}
\toprule
  \multirow{2}{*}{\bf Service}
    & {\bf Free space}
  & \textbf{Cost} \\
    & {\bf (GB)}
    & {\bf (\$/GB/year)}\\
\midrule
  Amazon S3
    & -
    & \$0.36 \\
  Box
    & \phantom{00}10 GB
    & \$0.60 \\
  Baidu
    & 2048 GB
    & \$0.80 \\
  Dropbox
    & \phantom{000}2 GB
    & \$1.20 \\
  Google Drive
    & \phantom{00}15 GB
    & \$0.24 \\
  OneDrive
    & \phantom{000}7 GB
    & \$0.50 \\
\bottomrule
\end{tabular}
\caption{A simple table.}
\label{t:simple}
\end{table}

\begin{table*}[t]
\centering
\footnotesize
\centering
\small
\setlength{\tabcolsep}{3pt}
\begin{tabular}{@{~}l l rrr c rrrr c rrrr@{~}}
\toprule
\multirow{2}{*}{Attack workload} &
\multirow{2}{*}{Trace} &
\multicolumn{3}{c}{Number of requests} & &
\multicolumn{4}{c}{Running time (seconds)} & &
\multicolumn{4}{c}{Number of data items}
\\ \cmidrule{3-5} \cmidrule{7-10} \cmidrule{12-15}
& & total & attacker & replayed &
  & original & replay & init. & loop &
          & accessed & reported 
          & false & missed
\\ \midrule
ACL error       & Submit.short
    & 1,746 & 45 & 17 &
    & 154.0 & 3.3 & 2.4  & 0.9 &
            & 428 & 193 & 0 & 0 \\
                & Submit.long
    & 7,071 & 45 & 17 &
    & 668.0 & 9.7 & 7.7 & 2.0 &
            & 428 & 193 & 0 & 0 \\
Stolen password & EndoApp.short
    & 1,219 & 25 & 110 &
    & 150.0 & 6.5 & 5.1 & 1.4 &
            & 596 & 137 & 1 & 0 \\
                & EndoApp.long
    & 6,193 & 25 & 219 &
    & 641.0 & 36.2 & 21.9 & 14.3 &
            & 2,472 & 197 & 1 & 0 \\
\bottomrule
\end{tabular}


\caption{A big table.}
\label{t:big}
\end{table*}
